\input format.tex

%\usepackage{arydshln}
\usepackage{graphicx}
\graphicspath{{cores/}}

\begin{document}
\color{gray2}

\vspace*{0mm}
\fontsize{9.3pt}{9.8pt}\selectfont
\begin{center}
{\bf\sanhao 健康建议}
\end{center}

\begin{tctabularx}{tabularx={m{15.4cm}}}
\\[-6.769pt]
\cellcolor{kelly} \raisebox{1pt}{\color{black70}\xiaosihao\lantxh\bf 肠道健康总说明} \\
\end{tctabularx}

{
\lantxh
\vspace*{-4.25mm}
\fontsize{8.8pt}{11pt}\selectfont
\begin{longtable}{|@{~~~}m{4cm}<{\centering}|m{10.96cm}|}
\hline
\parbox[c]{\hsize}{\vskip6pt\raisebox{-4ex}{\includegraphics[width=1.8cm]{INSUFFICIENCY-red.pdf}}~~\xiaowuhao \color{gray2} 消化和吸收\vskip6pt} &
\parbox[c]{\hsize}{\vskip6pt\xiaowuhao\color{gray2} {
通过菌群测得您的
蛋白质消化吸收能力偏低、脂肪消化吸收能力偏低、碳水化合物消化吸收能力偏低；15种主要代谢物中
蛋白质总量、谷胱甘肽代谢、牛磺酸代谢等指标存在异常。
}\vskip4pt}\\
\hline
\parbox[c]{\hsize}{\vskip6pt\raisebox{-4ex}{\includegraphics[width=1.8cm]{INFLAMMATION-red.pdf}}~~\xiaowuhao\color{gray2} 炎症和免疫\vskip6pt} &
\parbox[c]{\hsize}{\vskip6pt\xiaowuhao\color{gray2} {
您的肠道屏障功能、维持肠道屏障功能的能力、机体免疫力均较差。
}\vskip4pt}\\
\hline
\parbox[c]{\hsize}{\vskip6pt\raisebox{-4ex}{\includegraphics[width=1.8cm]{MICROBIOTA-low.pdf}}~~\xiaowuhao\color{gray2} 肠道菌群\vskip6pt} &
\parbox[c]{\hsize}{\vskip6pt\xiaowuhao\color{gray2} {
您的肠道菌群以普雷沃氏菌属为主。肠道菌群多样性水平
较差，菌群失调的风险较高。
有28类细菌指标与参考人群相比存在异常。
}\vskip4pt}\\
\hline
\parbox[c]{\hsize}{\vskip6pt\raisebox{-4ex}{\includegraphics[width=1.8cm]{INFECTION-yellow.pdf}}~~\xiaowuhao\color{gray2} 致病菌\vskip6pt} &
\parbox[c]{\hsize}{\vskip6pt\xiaowuhao\color{gray2} {
%您的肠道中检测到2种可能导致疾病的致病菌，其中种食源性致病菌，种机会致病菌，含量未超标，但在您免疫能力低下时仍可能会损害您的肠道健康。
您的肠道中检测到2种可能导致疾病的致病菌，含量未超标，但在您免疫能力低下时仍可能会损害您的肠道健康。
}\vskip4pt}\\
\hline
\parbox[c]{\hsize}{\vskip6pt\raisebox{-4ex}{\includegraphics[width=1.8cm]{IMBALANCE-red.pdf}}~~\xiaowuhao\color{gray2} 代谢平衡\vskip6pt} &
\parbox[c]{\hsize}{\vskip6pt\xiaowuhao\color{gray2} {
您的代谢失衡风险高，11种重要有机小分子代谢中维生素B1、维生素B2、维生素B3等代谢异常。
}\vskip4pt}\\
\hline
\end{longtable}
}

\fontsize{9.3pt}{9.8pt}\selectfont
\begin{tctabularx}{tabularx={m{15.4cm}}}
\\[-6.769pt]
\cellcolor{kelly} \raisebox{1pt}{\color{black70}\xiaosihao\lantxh\bf GIhealth 提示您} \\
\end{tctabularx}

{
\lantxh
\vspace*{-4.25mm}
\fontsize{8.8pt}{11pt}\selectfont
\begin{longtable}{|m{15.45cm}|}
\hline
\parbox[c]{\hsize}{\vskip11pt\begin{compactenum}[\mbox{1、}]
{\xiaowuhao\color{gray2}\item[1、]
您的肠道菌群多样性
较差，仅高于52.47{\%}的参考人群。肠道有益菌含量
较低，仅高于77.98{\%}的参考人群，有害菌含量
较低，高于26.56{\%}的参考人群。
检出2种致病菌，含量在正常范围内。
}
{\xiaowuhao\color{gray2}\item[2、]
通过对您的肠道菌群分析，您患2型糖尿病、结直肠癌、炎症性肠病的风险与参考人群一致。
}
{\xiaowuhao\color{gray2}\item[3、]
您肠道内参与维生素B1、维生素B2、维生素B3等代谢的菌群指标与参考人群相比存在异常。代谢失衡风险较高，需引起注意，建议您持续监测肠道健康。
}
{\xiaowuhao\color{gray2}\item[4、]
您肠道内参与蛋白质总量、谷胱甘肽代谢、牛磺酸代谢等的菌群指标与参考人群相比存在异常，需引起注意，建议您持续监测肠道健康。
}
{\xiaowuhao\color{gray2}\item[5、]}
\end{compactenum}\vskip2pt}\\
\hline
\end{longtable}
}

\end{document}
